% Generated by evaluation/exports.py - do not edit by hand.
% Source run:  130726_224804_agentic, 140726_132208_agentic, 140726_144552_agentic, 140726_171539_agentic, 140726_194805_agentic, 140726_215821_agentic
% Source data: samples.jsonl
% Command:     python -m evaluation.rebuild --all
% Rows:        600 samples
% Generated:   2026-08-01 18:03:09
\begin{table}[htbp]
  \centering
  \caption{Adjacent Pair Accuracy by organism, conditioned on fragment count. Mean protein length is shown per cell, since an organism gap that disappears within a fragment-count bin is a protein-length effect, not an organism effect.}
  \label{tab:organism_gap}
  \begin{tabular}{lrr}
    \toprule
    Fragments & E. coli & Yeast \\
    \midrule
    2-4 & 1.0000 (n=14, len 53) & 1.0000 (n=9, len 93) \\
    5-9 & 0.7915 (n=58, len 117) & 0.7543 (n=43, len 138) \\
    10-19 & 0.6593 (n=125, len 267) & 0.6701 (n=62, len 198) \\
    20-49 & 0.6010 (n=82, len 412) & 0.5793 (n=127, len 456) \\
    50+ & 0.6480 (n=20, len 861) & 0.4858 (n=58, len 1026) \\
    \bottomrule
  \end{tabular}
  \par\smallskip\footnotesize{Cells pool every replica count. The design is balanced - each organism contributes the same replica counts with equal sample counts - so the comparison is fair, but the pooling widens within-cell spread and these values are not comparable to a single run's numbers.}
\end{table}
